\documentclass[11pt]{article}
\usepackage[margin=1in]{geometry}
\usepackage{xcolor}
\usepackage{parskip}
\usepackage{fancyhdr}
\usepackage{amsmath}
\usepackage{needspace}

\pagestyle{fancy}
\fancyhf{}
\cfoot{\thepage}
\rhead{Lecture 1}
\lhead{MA 214: Applied Statistics}
\renewcommand{\headrulewidth}{.4mm}

% A key concept: bold heading, then open space to write, then a faint divider.
% \needspace keeps the heading + its writing space together: if less than ~1.5in
% remains on the page, the whole concept moves to the next page.
\newcommand{\concept}[1]{%
  \par\needspace{1.5in}\medskip\noindent{\large\bfseries #1}\par
  \vspace{1.15in}%
  \noindent{\color{black!15}\rule{\linewidth}{0.4pt}}\par\smallskip}

\begin{document}

\noindent\textbf{Name:}\underline{\hspace{2.25in}}\hfill\textbf{Date:}\underline{\hspace{1.4in}}

\begin{center}
 \Large Lecture 1 Notes: Randomness, Data, and Statistical Models
\end{center}

\noindent\textit{Room to record each key idea from today's lecture in your own words --- a definition,
an example, and any cautions. These pages are yours to keep.}
\medskip

\concept{Random processes}
\concept{Outcomes and events}
\concept{Probability (frequentist and Bayesian views)}
\concept{Basic probability rules}
\concept{Law of large numbers (proportions and means)}
\concept{Data-generating processes and statistical models}
\concept{Goals of statistical inference}
\concept{Observational studies vs.\ experiments}
\concept{Parameter, statistic, and point estimate}
\concept{Sampling variability and the sampling distribution}
\end{document}
